% Keep the closing discussion together after the completed experimental figures/tables.
\clearpage
\section{Discussion}
\label{sec:discussion}

The value of this characterization lies in distinguishing algorithmic
mean-shift contributions that the shared leading-order term leaves
unresolved. The explicit coefficients quantify the additional effect
of persistent controls, while the local-moment analysis explains how
it arises despite direct cancellation at the server.

For either algorithm's invariant law $\pi$ and $b=\E_\pi[x]-x^\star$,
parameter error decomposes as
\begin{equation}
\label{eq:stationary-bias-variance}
\E_\pi\|x-x^\star\|^2
=
\underbrace{\|b\|^2}_{\text{squared mean bias}}
+
\underbrace{\operatorname{tr}\Cov_\pi(x)}_{\text{stationary variance}}.
\end{equation}
Averaging stationary iterates preserves $b$, while temporal dependence governs
the variance of the average. The error share of a systematic displacement
therefore depends on how the final model is formed. The real-data studies
complement this characterization through finite-tail error decomposition
and post-hoc directional agreement with the derived curvature--noise
structure. They provide empirical interpretation, rather than identification
of the exact additive-noise coefficients under state-dependent minibatch
noise or certification of stationarity.

The formulas make the roles of client count, step size, local training
length, objective geometry, and noise explicit at the resolved orders.
This is an interpretable account of existing algorithms, not an overall
performance ranking: parameter MSE also depends on fluctuations and
averaging, while predictive loss depends on the learning problem.
